Enguany fa 40 anys de l'accident nuclear més greu de la història. La central nuclear de Txernòbil va alliberar una gran quantitat de materials radioactius a l'atmosfera. En particular, cesi-137 (${}^{137}_{\phantom{1}55}\mathrm{Cs}$), que en desintegrar-se mitjançant una emissió $\beta^-$ genera un nucli de bari. Aquest isòtop té una vida mitjana llarga (30,17 anys) i, per tant, representa un perill durant dècades.

\begin{apartat}{1,25}
Escriviu la reacció de desintegració del cesi-137. Calculeu quin percentatge de la quantitat inicial de cesi-137 queda avui, 40 anys després de l'accident.
\end{apartat}

\begin{apartat}{1,25}
En una zona a uns 10,0~km de la central es va detectar una concentració inicial de $6,00\times 10^{5}\ \mathrm{Bq\,m^{-2}}$ de cesi-137, immediatament després de l'accident. Segons alguns estudis, es considera segur tornar a habitar una zona afectada quan la concentració estigui per sota de $3,70\times 10^{4}\ \mathrm{Bq\,m^{-2}}$. Quin any podrà tornar-se a habitar aquesta zona si només considerem la desintegració natural del cesi-137?
\end{apartat}
